\begin{table}[htbp]
\centering
\caption{Share of agricultural heat exposure covered at a budget of 37 countries}
\label{tab:coverage}
\small
\begin{tabular}{lrrrr}
\toprule
 & & \multicolumn{3}{c}{Share of sample exposure (\%)}\\
\cmidrule(lr){3-5}
Rule & EMDEs & $\epsilon=0$ & $\epsilon=1$ & $\epsilon=1$\\
 & selected & & Bottom-first & Uniform\\
\midrule
Research-first          &  5 &  3.3 &  1.5 &  0.9 \\
Poorest-first           & 37 & 58.8 & 72.8 & 77.6 \\
Most-agricultural-first & 37 & 63.3 & 78.4 & 83.1 \\
Random (expected)       & 25.3 & 29.4 & 29.4 & 29.4 \\
\bottomrule
\end{tabular}
\begin{minipage}{0.95\textwidth}\footnotesize\smallskip
\emph{Notes.} The sample contains 126 countries, including 86 emerging and developing
economies. The rules rank countries by the country mean of log one
plus lagged climate-AI stock per million people, the country mean of lagged log GDP per capita,
or average agricultural employment share. Columns report
Equation~\eqref{eq:coverage}. At $\epsilon=0$, exposure is proportional to 2023 population
times agricultural employment share times the standard deviation of hottest-month
temperature anomalies. At $\epsilon=1$, the exposure assigned to each income decile
is weighted by the inverse of mean decile income. Bottom-first assigns agricultural
exposure to poorer deciles first; uniform allocates the same exposure fraction to
each decile. The shares use a common slope improvement and can be reproduced with
\texttt{code/31\_application\_agri.py --exposure-shares 0.3}.
\end{minipage}
\end{table}
